\newpage
\pagestyle{empty}

\section*{\centering РЕФЕРАТ}
\onehalfspacing
На 52 с., 11 рисунков, 3 таблицы, 0 приложений

Ключевые слова: методы Рунге--Кутты, законы сохранения, параметрические схемы, оптимизация, детерминированный сеточный поиск, метод роя частиц, нелинейное уравнение Шрёдингера, задача двух тел, задача~$N$~тел.

Тема выпускной квалификационной работы: «Оптимизация параметров метода Рунге-Кутты эволюционным алгоритмом».

В данной работе рассмотрен адаптивный параметрический метод Рунге--Кутты 6-го порядка и обоснована целесообразность его применения при решении задач с осциллирующими решениями. Исследованы два способа выбора наилучшего вектора свободных параметров: алгоритм детерминированного сеточного поиска и стохастический алгоритм роя частиц, а также предложены модификации и дополнения методов. Продемонстрированы численные решения четырёх физических задач, моделирующих осцилляционные процессы, проведено исследование точности разработанных методов и сравнение их эффективности как между собой, так и с непараметрическими методами Рунге--Кутты. На основе полученных численных решений проведены дополнительные эксперименты для оценки свойств методов и предложены направления для дальнейших исследований и развития подхода.

\newpage
\pagestyle{empty}

\section*{\centering ABSTRACT}
\onehalfspacing

52 pages, 11 figures, 3 tables, 0 appendices

Keywords: Runge--Kutta methods, conservation laws, parametric schemes, optimization, deterministic grid search, particle swarm method, nonlinear Schrödinger equation, two-body problem, $N$-body problem.

Master Thesis Topic: «Optimization of the Runge-Kutta method parameters using an evolutionary algorithm».

This paper examines an adaptive sixth-order parametric Runge-Kutta method and substantiates its applicability to problems with oscillatory solutions. Two approaches for selecting the optimal vector of free parameters are investigated: a deterministic grid search algorithm and a stochastic particle swarm optimization (PSO) algorithm. Furthermore, modifications and enhancements to these methods are proposed. Numerical solutions to four physical problems modeling oscillatory processes are presented, the accuracy of the developed methods is analyzed, and their efficiency is compared both with each other and with classical non-parametric Runge-Kutta methods. Based on the obtained numerical results, additional experiments are performed to evaluate the method properties, and directions for further research are outlined.